% GENERATED FILE. Edit canonical lessons, then run npm run generate:books.
% canonical-source-hash: fnv1a64:cf2d73fee70f28e0
% canonical-lessons: ES-C527-chocar, ES-C527-choque, ES-C527-atropellar, ES-C527-ileso, ES-C527-caer, ES-R527-repaso-520-527

\chapter{En la carretera}
\label{ch:es-en-la-carretera}
\hlchaptermodality{527}

\begin{chapteropening}
\textbf{By the end of this chapter:} \emph{I can describe a road accident: two cars crashed, someone was run over or fell, and who came out unhurt.}

Complete the last lesson of chapter 527: I can describe a road accident: two cars crashed, someone was run over or fell, and who came out unhurt.
\end{chapteropening}

\section[chocar]{chocar — to crash, to collide}

\label{lesson:ES-C527-chocar}



\begin{quote}
Before the new one: say the Spanish for to survive, then the Spanish for to trap.
\end{quote}

\subsection*{You'll want to know: chocar}
\textbf{chocar} — \textquotedblleft{}to crash, to collide\textquotedblright{}.

\begin{quote}
Dos coches chocaron en la esquina.
\end{quote}

\begin{quote}
\emph{Two cars crashed on the corner.}
\end{quote}

\emph{Chocar con} or \emph{chocar contra} — to crash into. \emph{¡Chócala!} — high five!

\begin{cousinweb}
The origin of \emph{chocar} is not settled. It is usually treated as imitative, a word that copies the noise of a blow; French \emph{choquer} is close beside it.
\end{cousinweb}

\begin{grammarlens}[title={What you've built}]
One more word for telling, explaining and describing at B1.
\end{grammarlens}

\subsection*{Your turn}
\begin{itemize}
\raggedright
  \item \emph{Say it:} \emph{chocar}
  \item \emph{Say it:} \emph{chocar}, once more
  \item \emph{Say it:} a sentence of your own with \emph{chocar}
  \item \emph{Recall:} say the Spanish for to survive, then the Spanish for to trap, then say \emph{chocar} again
\end{itemize}

\begin{tcolorbox}[breakable,colback=teal!4,colframe=teal!35!black,title={Before you move on}]
What does chocar mean? (\textquotedblleft{}To crash, to collide\textquotedblright{}.) Say it once more.
\end{tcolorbox}
\section[el choque]{el choque — the crash, the collision; the shock}

\label{lesson:ES-C527-choque}



\begin{quote}
Before the new one: say the Spanish for to trap, then the Spanish for to crash.
\end{quote}

\subsection*{You'll want to know: el choque}
\textbf{el choque} — \textquotedblleft{}the crash, the collision; the shock\textquotedblright{}.

\begin{quote}
El choque fue muy fuerte.
\end{quote}

\begin{quote}
\emph{The crash was very violent.}
\end{quote}

\emph{Un choque cultural} — a culture shock.

\begin{cousinweb}
\emph{choque} is built on the verb \emph{chocar} (EtymDB), whose own origin is not settled; it is often treated as a sound word.
\end{cousinweb}

\begin{grammarlens}[title={What you've built}]
One more word for telling, explaining and describing at B1.
\end{grammarlens}

\subsection*{Your turn}
\begin{itemize}
\raggedright
  \item \emph{Say it:} \emph{el choque}
  \item \emph{Say it:} \emph{el choque}, once more
  \item \emph{Say it:} a sentence of your own with \emph{el choque}
  \item \emph{Recall:} say the Spanish for to trap, then the Spanish for to crash, then say \emph{el choque} again
\end{itemize}

\begin{tcolorbox}[breakable,colback=teal!4,colframe=teal!35!black,title={Before you move on}]
What does el choque mean? (\textquotedblleft{}The crash, the collision; the shock\textquotedblright{}.) Say it once more.
\end{tcolorbox}
\section[atropellar]{atropellar — to run over, to knock down}

\label{lesson:ES-C527-atropellar}



\begin{quote}
Before the new one: say the Spanish for to crash, then the Spanish for the crash.
\end{quote}

\subsection*{You'll want to know: atropellar}
\textbf{atropellar} — \textquotedblleft{}to run over, to knock down\textquotedblright{}.

\begin{quote}
Un coche atropelló a una señora en el paso de cebra.
\end{quote}

\begin{quote}
\emph{A car knocked down a woman on the zebra crossing.}
\end{quote}

\emph{Me atropelló un coche} — I was hit by a car.

\begin{cousinweb}
\emph{atropellar} is built on \emph{tropel} 'a rushing crowd' (EtymDB), with \emph{a-}. \emph{Tropel} is usually said to have come into Spanish from Occitan or French, and its root is usually traced to Germanic.
\end{cousinweb}

\begin{grammarlens}[title={What you've built}]
One more word for telling, explaining and describing at B1.
\end{grammarlens}

\subsection*{Your turn}
\begin{itemize}
\raggedright
  \item \emph{Say it:} \emph{atropellar}
  \item \emph{Say it:} \emph{atropellar}, once more
  \item \emph{Say it:} a sentence of your own with \emph{atropellar}
  \item \emph{Recall:} say the Spanish for to crash, then the Spanish for the crash, then say \emph{atropellar} again
\end{itemize}

\begin{tcolorbox}[breakable,colback=teal!4,colframe=teal!35!black,title={Before you move on}]
What does atropellar mean? (\textquotedblleft{}To run over, to knock down\textquotedblright{}.) Say it once more.
\end{tcolorbox}
\section[ileso]{ileso — unhurt, unharmed}

\label{lesson:ES-C527-ileso}



\begin{quote}
Before the new one: say the Spanish for the crash, then the Spanish for to run over.
\end{quote}

\subsection*{You'll want to know: ileso}
\textbf{ileso} — \textquotedblleft{}unhurt, unharmed\textquotedblright{}.

\begin{quote}
Los pasajeros salieron ilesos del accidente.
\end{quote}

\begin{quote}
\emph{The passengers came out of the accident unhurt.}
\end{quote}

\emph{Salir ileso} — to escape unhurt.

\begin{cousinweb}
\emph{ileso} is a learned borrowing from Latin \emph{illaesus} 'unhurt, unharmed' (EtymDB), built on \emph{in-} 'not' and \emph{laesus} 'hurt', from \emph{laedere} 'to hurt'.
\end{cousinweb}

\begin{grammarlens}[title={What you've built}]
One more word for telling, explaining and describing at B1.
\end{grammarlens}

\subsection*{Your turn}
\begin{itemize}
\raggedright
  \item \emph{Say it:} \emph{ileso}
  \item \emph{Say it:} \emph{ileso}, once more
  \item \emph{Say it:} a sentence of your own with \emph{ileso}
  \item \emph{Recall:} say the Spanish for the crash, then the Spanish for to run over, then say \emph{ileso} again
\end{itemize}

\begin{tcolorbox}[breakable,colback=teal!4,colframe=teal!35!black,title={Before you move on}]
What does ileso mean? (\textquotedblleft{}Unhurt, unharmed\textquotedblright{}.) Say it once more.
\end{tcolorbox}
\section[caer]{caer — to fall}

\label{lesson:ES-C527-caer}



\begin{quote}
Before the new one: say the Spanish for to run over, then the Spanish for unhurt.
\end{quote}

\subsection*{You'll want to know: caer}
\textbf{caer} — \textquotedblleft{}to fall\textquotedblright{}.

\begin{quote}
Un árbol cayó en medio de la calle.
\end{quote}

\begin{quote}
\emph{A tree fell in the middle of the street.}
\end{quote}

\emph{Caerse} is to fall over: \emph{Me caí en la calle.} Forms: \emph{caigo}, \emph{cayó}, \emph{cayeron}.

\begin{cousinweb}
\emph{caer} is inherited from Latin \emph{cadere} 'to fall, to happen', through spoken Latin \emph{*cadēre} (EtymDB). The \emph{d} between vowels was lost.
\end{cousinweb}

\begin{grammarlens}[title={What you've built}]
One more word for telling, explaining and describing at B1.
\end{grammarlens}

\subsection*{Your turn}
\begin{itemize}
\raggedright
  \item \emph{Say it:} \emph{caer}
  \item \emph{Say it:} \emph{caer}, once more
  \item \emph{Say it:} a sentence of your own with \emph{caer}
  \item \emph{Recall:} say the Spanish for to run over, then the Spanish for unhurt, then say \emph{caer} again
\end{itemize}

\begin{tcolorbox}[breakable,colback=teal!4,colframe=teal!35!black,title={Before you move on}]
What does caer mean? (\textquotedblleft{}To fall\textquotedblright{}.) Say it once more.
\end{tcolorbox}
\section[(dialogue)]{Repaso de los capítulos 520–527}

\label{lesson:ES-R527-repaso-520-527}



\begin{quote}
Say the Spanish for the offence and the crime.
\end{quote}

\subsection*{Your turn}
Say these aloud:

\begin{itemize}
\raggedright
  \item the offence, the crime, the offender, the victim, the investigation
  \item to investigate, the clue, the proof, to find out, to suspect
  \item the trial, the judge, to judge, to sentence, the prison
  \item the prisoner, to accuse, to confess, to testify, to free
  \item the fire, to set on fire, to burn, the alarm, to evacuate
  \item the wound, to wound, the ambulance, to help, the injured person
  \item to rescue, the rescue, to save, to survive, to trap
  \item to crash, the crash, to run over, unhurt, to fall
\end{itemize}

\begin{tcolorbox}[breakable,colback=teal!4,colframe=teal!35!black,title={Before you move on}]
The offence? (\textbf{el delito}.) The crime? (\textbf{el crimen}.) The offender? (\textbf{el delincuente}.) The victim? (\textbf{la víctima}.) The investigation? (\textbf{la investigación}.) To investigate? (\textbf{investigar}.) The clue? (\textbf{la pista}.) The proof? (\textbf{la prueba}.) To find out? (\textbf{averiguar}.) To suspect? (\textbf{sospechar}.) The trial? (\textbf{el juicio}.) The judge? (\textbf{el juez}.) To judge? (\textbf{juzgar}.) To sentence? (\textbf{condenar}.) The prison? (\textbf{la cárcel}.) The prisoner? (\textbf{el preso}.) To accuse? (\textbf{acusar}.) To confess? (\textbf{confesar}.) To testify? (\textbf{testificar}.) To free? (\textbf{liberar}.) The fire? (\textbf{el incendio}.) To set on fire? (\textbf{incendiar}.) To burn? (\textbf{arder}.) The alarm? (\textbf{la alarma}.) To evacuate? (\textbf{evacuar}.) The wound? (\textbf{la herida}.) To wound? (\textbf{herir}.) The ambulance? (\textbf{la ambulancia}.) To help? (\textbf{socorrer}.) The injured person? (\textbf{el herido}.) To rescue? (\textbf{rescatar}.) The rescue? (\textbf{el rescate}.) To save? (\textbf{salvar}.) To survive? (\textbf{sobrevivir}.) To trap? (\textbf{atrapar}.) To crash? (\textbf{chocar}.) The crash? (\textbf{el choque}.) To run over? (\textbf{atropellar}.) Unhurt? (\textbf{ileso}.) To fall? (\textbf{caer}.)
\end{tcolorbox}
